% Options for packages loaded elsewhere
\PassOptionsToPackage{unicode}{hyperref}
\PassOptionsToPackage{hyphens}{url}
\documentclass[
]{article}
\usepackage{xcolor}
\usepackage[margin=1in]{geometry}
\usepackage{amsmath,amssymb}
\setcounter{secnumdepth}{-\maxdimen} % remove section numbering
\usepackage{iftex}
\ifPDFTeX
  \usepackage[T1]{fontenc}
  \usepackage[utf8]{inputenc}
  \usepackage{textcomp} % provide euro and other symbols
\else % if luatex or xetex
  \usepackage{unicode-math} % this also loads fontspec
  \defaultfontfeatures{Scale=MatchLowercase}
  \defaultfontfeatures[\rmfamily]{Ligatures=TeX,Scale=1}
\fi
\usepackage{lmodern}
\ifPDFTeX\else
  % xetex/luatex font selection
\fi
% Use upquote if available, for straight quotes in verbatim environments
\IfFileExists{upquote.sty}{\usepackage{upquote}}{}
\IfFileExists{microtype.sty}{% use microtype if available
  \usepackage[]{microtype}
  \UseMicrotypeSet[protrusion]{basicmath} % disable protrusion for tt fonts
}{}
\makeatletter
\@ifundefined{KOMAClassName}{% if non-KOMA class
  \IfFileExists{parskip.sty}{%
    \usepackage{parskip}
  }{% else
    \setlength{\parindent}{0pt}
    \setlength{\parskip}{6pt plus 2pt minus 1pt}}
}{% if KOMA class
  \KOMAoptions{parskip=half}}
\makeatother
\usepackage{color}
\usepackage{fancyvrb}
\newcommand{\VerbBar}{|}
\newcommand{\VERB}{\Verb[commandchars=\\\{\}]}
\DefineVerbatimEnvironment{Highlighting}{Verbatim}{commandchars=\\\{\}}
% Add ',fontsize=\small' for more characters per line
\usepackage{framed}
\definecolor{shadecolor}{RGB}{248,248,248}
\newenvironment{Shaded}{\begin{snugshade}}{\end{snugshade}}
\newcommand{\AlertTok}[1]{\textcolor[rgb]{0.94,0.16,0.16}{#1}}
\newcommand{\AnnotationTok}[1]{\textcolor[rgb]{0.56,0.35,0.01}{\textbf{\textit{#1}}}}
\newcommand{\AttributeTok}[1]{\textcolor[rgb]{0.13,0.29,0.53}{#1}}
\newcommand{\BaseNTok}[1]{\textcolor[rgb]{0.00,0.00,0.81}{#1}}
\newcommand{\BuiltInTok}[1]{#1}
\newcommand{\CharTok}[1]{\textcolor[rgb]{0.31,0.60,0.02}{#1}}
\newcommand{\CommentTok}[1]{\textcolor[rgb]{0.56,0.35,0.01}{\textit{#1}}}
\newcommand{\CommentVarTok}[1]{\textcolor[rgb]{0.56,0.35,0.01}{\textbf{\textit{#1}}}}
\newcommand{\ConstantTok}[1]{\textcolor[rgb]{0.56,0.35,0.01}{#1}}
\newcommand{\ControlFlowTok}[1]{\textcolor[rgb]{0.13,0.29,0.53}{\textbf{#1}}}
\newcommand{\DataTypeTok}[1]{\textcolor[rgb]{0.13,0.29,0.53}{#1}}
\newcommand{\DecValTok}[1]{\textcolor[rgb]{0.00,0.00,0.81}{#1}}
\newcommand{\DocumentationTok}[1]{\textcolor[rgb]{0.56,0.35,0.01}{\textbf{\textit{#1}}}}
\newcommand{\ErrorTok}[1]{\textcolor[rgb]{0.64,0.00,0.00}{\textbf{#1}}}
\newcommand{\ExtensionTok}[1]{#1}
\newcommand{\FloatTok}[1]{\textcolor[rgb]{0.00,0.00,0.81}{#1}}
\newcommand{\FunctionTok}[1]{\textcolor[rgb]{0.13,0.29,0.53}{\textbf{#1}}}
\newcommand{\ImportTok}[1]{#1}
\newcommand{\InformationTok}[1]{\textcolor[rgb]{0.56,0.35,0.01}{\textbf{\textit{#1}}}}
\newcommand{\KeywordTok}[1]{\textcolor[rgb]{0.13,0.29,0.53}{\textbf{#1}}}
\newcommand{\NormalTok}[1]{#1}
\newcommand{\OperatorTok}[1]{\textcolor[rgb]{0.81,0.36,0.00}{\textbf{#1}}}
\newcommand{\OtherTok}[1]{\textcolor[rgb]{0.56,0.35,0.01}{#1}}
\newcommand{\PreprocessorTok}[1]{\textcolor[rgb]{0.56,0.35,0.01}{\textit{#1}}}
\newcommand{\RegionMarkerTok}[1]{#1}
\newcommand{\SpecialCharTok}[1]{\textcolor[rgb]{0.81,0.36,0.00}{\textbf{#1}}}
\newcommand{\SpecialStringTok}[1]{\textcolor[rgb]{0.31,0.60,0.02}{#1}}
\newcommand{\StringTok}[1]{\textcolor[rgb]{0.31,0.60,0.02}{#1}}
\newcommand{\VariableTok}[1]{\textcolor[rgb]{0.00,0.00,0.00}{#1}}
\newcommand{\VerbatimStringTok}[1]{\textcolor[rgb]{0.31,0.60,0.02}{#1}}
\newcommand{\WarningTok}[1]{\textcolor[rgb]{0.56,0.35,0.01}{\textbf{\textit{#1}}}}
\usepackage{graphicx}
\makeatletter
\newsavebox\pandoc@box
\newcommand*\pandocbounded[1]{% scales image to fit in text height/width
  \sbox\pandoc@box{#1}%
  \Gscale@div\@tempa{\textheight}{\dimexpr\ht\pandoc@box+\dp\pandoc@box\relax}%
  \Gscale@div\@tempb{\linewidth}{\wd\pandoc@box}%
  \ifdim\@tempb\p@<\@tempa\p@\let\@tempa\@tempb\fi% select the smaller of both
  \ifdim\@tempa\p@<\p@\scalebox{\@tempa}{\usebox\pandoc@box}%
  \else\usebox{\pandoc@box}%
  \fi%
}
% Set default figure placement to htbp
\def\fps@figure{htbp}
\makeatother
\setlength{\emergencystretch}{3em} % prevent overfull lines
\providecommand{\tightlist}{%
  \setlength{\itemsep}{0pt}\setlength{\parskip}{0pt}}
\usepackage{bookmark}
\IfFileExists{xurl.sty}{\usepackage{xurl}}{} % add URL line breaks if available
\urlstyle{same}
\hypersetup{
  pdftitle={UMAP and t-SNE},
  hidelinks,
  pdfcreator={LaTeX via pandoc}}

\title{UMAP and t-SNE}
\author{}
\date{\vspace{-2.5em}}

\begin{document}
\maketitle

\textbf{Workflow labs} use the variant template: Goal → Approach →
Execution → Check → Report.

\subsection{Learning objectives}\label{learning-objectives}

\begin{itemize}
\tightlist
\item
  Compute UMAP and t-SNE embeddings and describe what each optimisation
  is doing.
\item
  Interpret global vs local structure in a UMAP/t-SNE plot and avoid
  over-reading cluster distances.
\item
  Tune \texttt{n\_neighbors} / \texttt{perplexity} and see the visual
  effect.
\end{itemize}

\subsection{Prerequisites}\label{prerequisites}

PCA; basic notion of neighbour graphs.

\subsection{Background}\label{background}

UMAP and t-SNE are nonlinear embedding methods widely used in
single-cell and imaging biomedicine. Both build a neighbour graph in
high dimensions and then try to match a lower-dimensional graph with
similar neighbour structure. They are not clustering algorithms, and
distances in the embedding are not Euclidean distances in the original
space. The prominent visual clusters in a UMAP often exist in the data,
but the distance between clusters is close to meaningless and the
density inside a cluster is not comparable across clusters.

t-SNE emphasises local structure; UMAP, with default settings, is a
little more faithful to global structure but can still distort. Both
have knobs --- \texttt{perplexity} for t-SNE, \texttt{n\_neighbors} and
\texttt{min\_dist} for UMAP --- that change the resulting picture
substantially. A reviewer who sees only one embedding has no way to tell
if the picture is a stable feature of the data or an artefact of the
settings.

Treat UMAP/t-SNE as a visualisation of neighbourhood structure, not as a
geometric truth. Do not run statistical tests on the coordinates. Do not
claim a trajectory unless you used a trajectory method.

\subsection{Setup}\label{setup}

\begin{Shaded}
\begin{Highlighting}[]
\FunctionTok{library}\NormalTok{(tidyverse)}
\FunctionTok{library}\NormalTok{(uwot)}
\FunctionTok{library}\NormalTok{(Rtsne)}
\FunctionTok{set.seed}\NormalTok{(}\DecValTok{42}\NormalTok{)}
\FunctionTok{theme\_set}\NormalTok{(}\FunctionTok{theme\_minimal}\NormalTok{(}\AttributeTok{base\_size =} \DecValTok{12}\NormalTok{))}
\end{Highlighting}
\end{Shaded}

\subsection{1. Goal}\label{goal}

Embed a simulated high-dimensional dataset with three blobs into two
dimensions using UMAP and t-SNE, and compare against PCA.

\subsection{2. Approach}\label{approach}

\begin{Shaded}
\begin{Highlighting}[]
  \FunctionTok{rep}\NormalTok{(mu, }\AttributeTok{each =}\NormalTok{ n)}
\NormalTok{X }\OtherTok{\textless{}{-}} \FunctionTok{rbind}\NormalTok{(}
  \FunctionTok{make\_blob}\NormalTok{(}\DecValTok{100}\NormalTok{, }\FunctionTok{c}\NormalTok{(}\DecValTok{0}\NormalTok{, }\DecValTok{0}\NormalTok{, }\FunctionTok{rep}\NormalTok{(}\DecValTok{0}\NormalTok{, }\DecValTok{18}\NormalTok{))),}
  \FunctionTok{make\_blob}\NormalTok{(}\DecValTok{100}\NormalTok{, }\FunctionTok{c}\NormalTok{(}\DecValTok{4}\NormalTok{, }\DecValTok{0}\NormalTok{, }\FunctionTok{rep}\NormalTok{(}\DecValTok{0}\NormalTok{, }\DecValTok{18}\NormalTok{))),}
  \FunctionTok{make\_blob}\NormalTok{(}\DecValTok{100}\NormalTok{, }\FunctionTok{c}\NormalTok{(}\DecValTok{0}\NormalTok{, }\DecValTok{4}\NormalTok{, }\FunctionTok{rep}\NormalTok{(}\DecValTok{0}\NormalTok{, }\DecValTok{18}\NormalTok{)))}
\NormalTok{)}
\NormalTok{lbl }\OtherTok{\textless{}{-}} \FunctionTok{rep}\NormalTok{(}\FunctionTok{c}\NormalTok{(}\StringTok{"A"}\NormalTok{, }\StringTok{"B"}\NormalTok{, }\StringTok{"C"}\NormalTok{), }\AttributeTok{each =} \DecValTok{100}\NormalTok{)}
\end{Highlighting}
\end{Shaded}

\subsection{3. Execution}\label{execution}

\begin{Shaded}
\begin{Highlighting}[]
\NormalTok{emb\_pca  }\OtherTok{\textless{}{-}} \FunctionTok{prcomp}\NormalTok{(X)}\SpecialCharTok{$}\NormalTok{x[, }\DecValTok{1}\SpecialCharTok{:}\DecValTok{2}\NormalTok{]}
\NormalTok{emb\_umap }\OtherTok{\textless{}{-}} \FunctionTok{umap}\NormalTok{(X, }\AttributeTok{n\_neighbors =} \DecValTok{15}\NormalTok{, }\AttributeTok{min\_dist =} \FloatTok{0.1}\NormalTok{)}
\NormalTok{emb\_tsne }\OtherTok{\textless{}{-}} \FunctionTok{Rtsne}\NormalTok{(X, }\AttributeTok{perplexity =} \DecValTok{30}\NormalTok{, }\AttributeTok{check\_duplicates =} \ConstantTok{FALSE}\NormalTok{)}\SpecialCharTok{$}\NormalTok{Y}

\NormalTok{df }\OtherTok{\textless{}{-}} \FunctionTok{bind\_rows}\NormalTok{(}
  \FunctionTok{tibble}\NormalTok{(}\AttributeTok{method =} \StringTok{"PCA"}\NormalTok{,  }\AttributeTok{x =}\NormalTok{ emb\_pca[, }\DecValTok{1}\NormalTok{],  }\AttributeTok{y =}\NormalTok{ emb\_pca[, }\DecValTok{2}\NormalTok{],  }\AttributeTok{lbl =}\NormalTok{ lbl),}
  \FunctionTok{tibble}\NormalTok{(}\AttributeTok{method =} \StringTok{"UMAP"}\NormalTok{, }\AttributeTok{x =}\NormalTok{ emb\_umap[, }\DecValTok{1}\NormalTok{], }\AttributeTok{y =}\NormalTok{ emb\_umap[, }\DecValTok{2}\NormalTok{], }\AttributeTok{lbl =}\NormalTok{ lbl),}
  \FunctionTok{tibble}\NormalTok{(}\AttributeTok{method =} \StringTok{"tSNE"}\NormalTok{, }\AttributeTok{x =}\NormalTok{ emb\_tsne[, }\DecValTok{1}\NormalTok{], }\AttributeTok{y =}\NormalTok{ emb\_tsne[, }\DecValTok{2}\NormalTok{], }\AttributeTok{lbl =}\NormalTok{ lbl)}
\NormalTok{)}
\FunctionTok{ggplot}\NormalTok{(df, }\FunctionTok{aes}\NormalTok{(x, y, }\AttributeTok{colour =}\NormalTok{ lbl)) }\SpecialCharTok{+}
  \FunctionTok{geom\_point}\NormalTok{(}\AttributeTok{alpha =} \FloatTok{0.7}\NormalTok{, }\AttributeTok{size =} \FloatTok{0.8}\NormalTok{) }\SpecialCharTok{+}
  \FunctionTok{facet\_wrap}\NormalTok{(}\SpecialCharTok{\textasciitilde{}}\NormalTok{ method, }\AttributeTok{scales =} \StringTok{"free"}\NormalTok{) }\SpecialCharTok{+}
  \FunctionTok{labs}\NormalTok{(}\AttributeTok{x =} \ConstantTok{NULL}\NormalTok{, }\AttributeTok{y =} \ConstantTok{NULL}\NormalTok{)}
\end{Highlighting}
\end{Shaded}

\subsection{4. Check}\label{check}

Sensitivity to \texttt{n\_neighbors}.

\begin{Shaded}
\begin{Highlighting}[]
\NormalTok{do\_umap }\OtherTok{\textless{}{-}} \ControlFlowTok{function}\NormalTok{(k) \{}
\NormalTok{  e }\OtherTok{\textless{}{-}} \FunctionTok{umap}\NormalTok{(X, }\AttributeTok{n\_neighbors =}\NormalTok{ k, }\AttributeTok{min\_dist =} \FloatTok{0.1}\NormalTok{)}
  \FunctionTok{tibble}\NormalTok{(}\AttributeTok{x =}\NormalTok{ e[, }\DecValTok{1}\NormalTok{], }\AttributeTok{y =}\NormalTok{ e[, }\DecValTok{2}\NormalTok{], }\AttributeTok{lbl =}\NormalTok{ lbl, }\AttributeTok{k =} \FunctionTok{paste0}\NormalTok{(}\StringTok{"nn="}\NormalTok{, k))}
\NormalTok{\}}
\FunctionTok{bind\_rows}\NormalTok{(}\FunctionTok{do\_umap}\NormalTok{(}\DecValTok{5}\NormalTok{), }\FunctionTok{do\_umap}\NormalTok{(}\DecValTok{15}\NormalTok{), }\FunctionTok{do\_umap}\NormalTok{(}\DecValTok{50}\NormalTok{)) }\SpecialCharTok{|\textgreater{}}
  \FunctionTok{ggplot}\NormalTok{(}\FunctionTok{aes}\NormalTok{(x, y, }\AttributeTok{colour =}\NormalTok{ lbl)) }\SpecialCharTok{+} \FunctionTok{geom\_point}\NormalTok{(}\AttributeTok{alpha =} \FloatTok{0.7}\NormalTok{, }\AttributeTok{size =} \FloatTok{0.8}\NormalTok{) }\SpecialCharTok{+}
  \FunctionTok{facet\_wrap}\NormalTok{(}\SpecialCharTok{\textasciitilde{}}\NormalTok{ k, }\AttributeTok{scales =} \StringTok{"free"}\NormalTok{)}
\end{Highlighting}
\end{Shaded}

\subsection{5. Report}\label{report}

\begin{quote}
Both UMAP and t-SNE recovered the three simulated blobs as distinct
visual clusters; PCA separated them more weakly because the signal lay
along two of 20 dimensions. UMAP embeddings with \texttt{n\_neighbors} ∈
\{5, 15, 50\} all identified the same three groups but differed
substantially in their relative layout.
\end{quote}

The message is not which method is best, but that the embedding is best
used as a visual companion, with another method --- clustering,
classification, or differential analysis --- doing the statistical work.

\subsection{Common pitfalls}\label{common-pitfalls}

\begin{itemize}
\tightlist
\item
  Reading distances between clusters off a UMAP as if they were
  biological distances.
\item
  Running one embedding with defaults and treating the result as the
  data.
\item
  Using UMAP coordinates as features in a downstream supervised model.
\end{itemize}

\subsection{Further reading}\label{further-reading}

\begin{itemize}
\tightlist
\item
  McInnes L, Healy J, Melville J (2018), \emph{UMAP: Uniform manifold
  approximation and projection}.
\item
  Kobak D, Berens P (2019), \emph{The art of using t-SNE for single-cell
  transcriptomics}.
\end{itemize}

\subsection{Session info}\label{session-info}

\begin{Shaded}
\begin{Highlighting}[]

\end{Highlighting}
\end{Shaded}

\subsection{See also --- chapter index}\label{see-also-chapter-index}

\begin{itemize}
\tightlist
\item
  \href{../index.Rmd}{Methods in this chapter}
\item
  \href{../../../ch00-find-your-method.Rmd}{Find your method (Chapter
  0)}
\end{itemize}

\end{document}
